\chapter{راهنمای اجرا و شواهد تکمیلی}
\thispagestyle{empty}

\section{اجرای محلی}
اجرای محلی پروژه با \lr{Docker Compose} انجام می‌شود:
\begin{latin}
\begin{lstlisting}
docker compose up --build
\end{lstlisting}
\end{latin}

در این پیکربندی، فرانت‌اند روی \lr{http://localhost:5173} و بک‌اند روی \lr{http://localhost:8000/api} منتشر می‌شود.

این روش برای بررسی سریع کل سامانه مناسب است، زیرا هر دو بخش اصلی را هم‌زمان بالا می‌آورد. پس از اجرا، کاربر می‌تواند ابتدا مسیر سلامت بک‌اند را بررسی کند و سپس از طریق فرانت‌اند وارد جریان مرور کاتالوگ، ورود، داشبورد و ابزارهای توسعه‌دهنده شود. اگر یکی از سرویس‌ها اجرا نشود، باید لاگ همان سرویس در \lr{Docker Compose} بررسی شود.

\section{اجرای آزمون بک‌اند}
دستور اجراشده در این گزارش:
\begin{latin}
\begin{lstlisting}
python manage.py test api
\end{lstlisting}
\end{latin}
خروجی خلاصه: \lr{Ran 49 tests} و \lr{OK}.

این دستور فقط آزمون‌های برنامه \lr{api} را اجرا می‌کند و برای تکرار نتیجه فصل آزمون کافی است. در صورت تغییر کد بک‌اند، اجرای دوباره این دستور باید نخستین کنترل کیفی باشد. برای تغییرات فرانت‌اند، اجرای آزمون‌های مرورگر و ساخت نسخه فرانت‌اند نیز باید به این کنترل اضافه شود.

\section{مسیرهای مهم}
\begin{longtable}{p{6cm} p{6cm}}
\caption{مسیرهای مهم برای بررسی دستی سامانه}\label{tab:runbook-routes}\\
\toprule
\rowcolor{tablehead}\textbf{مسیر} & \textbf{کاربرد} \\
\midrule
\ModernTableStart
\lr{/api/v1/system/health/} & سلامت سامانه \\
\lr{/api/v1/schema/openapi.json} & طرح \lr{OpenAPI} \\
\lr{/api/v1/auth/login/} & ورود نشست‌محور \\
\lr{/api/v1/catalog/apis/} & فهرست و انتشار \lr{API} \\
\lr{/api/v1/account/usage/} & تاریخچه مصرف \\
\lr{/api/v1/account/studio/flows/} & جریان‌های \lr{Studio} \\
\ModernTableStop
\bottomrule
\end{longtable}
